\qhead{11. Legal and Ethical Issues on Pricing — What are your company's main ethical challenges concerning pricing? How do they affect pricing practices? Does your company face legal cases, and do you see potential for such challenges?}

I evaluate three allegations using the objectionable-pricing taxonomy \parencite{nagle2018}.

\textbf{(a) Predatory pricing (``too low'').} The economic test is price below average variable cost \emph{plus} intent to recoup after rivals exit. BYD's prices are low but it stays \textbf{gross-margin positive} (about 17--20\%), so ``below cost'' is hard to establish — its structural cost edge means low price \(\neq\) below cost, and recoupment is unproven. Aggressive, but not legally predatory for BYD specifically.

\textbf{(b) Dumping / unfair subsidy.} This is the EU's actual finding: the anti-subsidy investigation imposed \textbf{countervailing duties of 17\% on BYD} — the lowest of the sampled firms, versus 35.3\% on SAIC — for five years from October 2024 \parencite{ec2024}. The objection is to the subsidised cost base, not BYD's own conduct, which directly affects practice by pushing BYD to localise production in Europe.

\textbf{(c) Ethical lens (six criteria, plus the firm-theory debate).} \emph{Utilitarian:} cheap EVs raise consumer welfare and cut emissions. \emph{Justice/fairness:} EU producers call subsidised pricing unfair and Chinese rivals call the cuts destructive — China's regulators publicly condemned ``involution'' (\emph{neijuan}) and warned against below-cost selling \parencite{cnbcInvolution}, even moving to reform the Pricing Law. \emph{Self-interest:} the 19\% profit fall shows even BYD's long-run interest is threatened \parencite{byd2025}. This maps onto Friedman's shareholder logic (maximise share within the law) versus Freeman's stakeholder view, which the state invokes to restrain BYD on industry-health grounds.

\textbf{Potential challenge.} If BYD became \emph{dominant} in an EU market (above the \(\sim\)40\% presumption threshold), \textbf{Article 102 TFEU} — abuse of dominance, covering predatory pricing and price discrimination — would apply, a sharper constraint than US law, which has no standalone offence for merely holding market power.
